\section*{अध्याय \ref{ch:b1:stereochemistry-in-depth} --- त्रिविम रसायन का गहन अध्ययन}

\begin{solution}{exo:b1:stereochemistry-in-depth:1}
\omterm{def:b1:stereochemistry-in-depth:isomers}{संरचनात्मक समावयवी}; \omterm{def:b1:stereochemistry-in-depth:enantiomers}{प्रतिबिंबरूपी}; \omterm{def:b1:stereochemistry-in-depth:enantiomers}{अप्रतिबिंबरूपी} (ऐसे \omterm{def:b1:stereochemistry-in-depth:isomers}{त्रिविम समावयवी} जो
दर्पण-प्रतिबिंब नहीं हैं); एक ही यौगिक के दो \omterm{def:b1:stereochemistry-in-depth:isomers}{संरूपण}; \omterm{def:b1:stereochemistry-in-depth:enantiomers}{अप्रतिबिंबरूपी}
(\cip{R,S} मेसो रूप है)।
\end{solution}

\begin{solution}{exo:b1:stereochemistry-in-depth:2}
\ce{-OH} $>$ \ce{-NH2} $>$ \ce{-CH3} $>$ \ce{-H} (O, N, C, H)।
\ce{-COOH} (O,O,O) $>$ \ce{-CHO} (O,O,H) $>$ \ce{-CH2OH} (O,H,H) $>$
\ce{-CH(CH3)2} (C,C,H)। \ce{-Br} $>$ \ce{-Cl} $>$ \ce{-CCl3} (Cl,Cl,Cl)
$>$ \ce{-CH2Cl} (Cl,H,H): पहला परमाणु अपने पड़ोसियों से पहले निर्णय करता है।
\end{solution}

\begin{solution}{exo:b1:stereochemistry-in-depth:3}
2,3-डाइक्लोरोब्यूटेन: दो तुल्य केंद्र, \cip{R,R}, \cip{S,S} और
मेसो \cip{R,S}: तीन। 2-ब्रोमो-3-क्लोरोब्यूटेन: दो भिन्न केंद्र,
चार \omterm{def:b1:stereochemistry-in-depth:isomers}{त्रिविम समावयवी}, कोई मेसो रूप नहीं। पेंटेन-2,4-डाइऑल: तीन, जिनमें \cip{2R,4S}
मेसो है।
\end{solution}

\begin{solution}{exo:b1:stereochemistry-in-depth:4}
C5 पर: आइसोप्रोपेनिल (C,C,C, द्वि-आबंध दो बार गिना जाता है) $>$ C6 $>$ C4
$>$ H। C6 और C4 दोनों पर (C,H,H) है; एक कदम आगे C6
कार्बोनिल कार्बन (O,O,C) तक ले जाता है और C4, C3 (C,C,H) तक: C6 जीतता है। आइसोप्रोपेनिल
समूह सामने और H पीछे होने पर, क्रम आइसोप्रोपेनिल $\to$ C6 $\to$ C4
जैसा खींचा गया है वैसा दक्षिणावर्त घूमता है: \cip{R}, पुदीने वाला कार्वोन; धारीदार फन्नी वाला
\cip{S} है।
\end{solution}

\begin{solution}{exo:b1:stereochemistry-in-depth:5}
$\theta = 0$: मेथिल ग्रसित, \qty{16.5}{kJ/mol}; $60^\circ$: गाउश, लगभग
\qty{2.9}{kJ/mol} ($62^\circ$ पर न्यूनतम 2.8); $120^\circ$: मेथिल द्वारा H का ग्रसन,
\qty{15.2}{kJ/mol}; $180^\circ$: प्रति, 0। ऊर्जाएँ बराबर होने पर, एक प्रति के लिए दो \omterm{def:b1:stereochemistry-in-depth:conformations}{गाउश
संरूपण}: दो-तिहाई गाउश।
\end{solution}

\begin{solution}{exo:b1:stereochemistry-in-depth:6}
शृंखला ऊर्ध्वाधर, \ce{COOH} शीर्ष पर और \ce{CH3} तल पर; क्रम
$\ce{NH2} > \ce{COOH} > \ce{CH3} > \ce{H}$। \ce{NH2} को बाईं ओर
और H को दाईं ओर रखने पर, $\ce{NH2} \to \ce{COOH} \to \ce{CH3}$ जैसा दिखता है वैसा
दक्षिणावर्त घूमता है; H क्षैतिज है (पाठक की ओर), इसलिए वर्णक
उलटता है: \cip{S}। ऐमीनो समूह बाईं ओर है।
\end{solution}

\begin{solution}{exo:b1:stereochemistry-in-depth:7}
\textit{cis}: सन्निकट कार्बनों पर \omterm{def:b1:stereochemistry-in-depth:conformations}{अक्षीय} और विषुवतीय स्थितियों में एक आबंध ऊपर और एक
नीचे होने से दोनों \omterm{def:b1:stereochemistry-in-depth:conformations}{कुर्सियों} में अक्षीय--विषुवतीय मिलता है।
\textit{trans}: द्विविषुवतीय या द्विअक्षीय। द्विविषुवतीय \omterm{def:b1:stereochemistry-in-depth:conformations}{कुर्सी} में दोनों
मेथिल एक-दूसरे के प्रति गाउश हैं (एक अन्योन्यक्रिया); \textit{cis}
समावयवी में वह अन्योन्यक्रिया और साथ में एक \omterm{def:b1:stereochemistry-in-depth:conformations}{अक्षीय} मेथिल है, लगभग \qty{5.7}{kJ/mol}
अधिक। \textit{trans} समावयवी अधिक स्थायी है।
\end{solution}

\begin{solution}{exo:b1:stereochemistry-in-depth:8}
$[\alpha] = 13.2/(2.00 \times 0.100) = +66.0^\circ$। \qty{1.00}{dm} की नली:
$+6.60^\circ$। \omterm{def:b1:stereochemistry-in-depth:enantiomers}{प्रतिबिंबरूपी}: \qty{2.00}{dm} की नली में $-13.2^\circ$।
\end{solution}

\begin{solution}{exo:b1:stereochemistry-in-depth:9}
\qty{25}{\celsius}: $\exp(5700/(8.314 \times 298.15)) = 10.0$, अंश
$10.0/11.0 = \qty{91}{\percent}$। \qty{-50}{\celsius}: $\exp(5700/(8.314
\times 223.15)) = 21.6$, \qty{96}{\percent}।
\end{solution}

\begin{solution}{exo:b1:stereochemistry-in-depth:10}
एथेन में सभी छह हाइड्रोजन एक जैसे हैं: हर \omterm{def:b1:stereochemistry-in-depth:conformations}{सांतरित संरूपण}
वही है। ब्यूटेन में पीछे वाला मेथिल सामने वाले के विपरीत खड़ा हो सकता है
(प्रति) या उसके बगल में ($\pm 60^\circ$, गाउश), जहाँ दोनों मेथिल एक-दूसरे पर
भीड़ करते हैं। जनसंख्याएँ: प्रति 1, गाउश $2\exp(-2.8/2.48) = 0.65$: लगभग
\qty{61}{\percent} प्रति और \qty{39}{\percent} गाउश।
\end{solution}

\begin{solution}{exo:b1:stereochemistry-in-depth:11}
C2 और C4 त्रिविमजनक हैं। C3 पर H, OH और दो शाखाएँ हैं जो
संरचना में सर्वसम हैं; वे केवल तभी भिन्न होती हैं जब C2 और C4 के वर्णक
विपरीत हों, और केवल तभी C3 त्रिविमजनक है (छद्म-असममित)। समावयवी:
\cip{2R,4R} और \cip{2S,4S}, \omterm{def:b1:stereochemistry-in-depth:enantiomers}{प्रतिबिंबरूपियों} का एक युग्म (C3 त्रिविमजनक नहीं);
\cip{2R,4S}, C3 पर दो व्यवस्थाओं के साथ, दो मेसो रूप। कुल चार।
\end{solution}

\begin{solution}{exo:b1:stereochemistry-in-depth:12}
$[\alpha]_{\text{नमूना}} = 2.31/(1.00 \times 0.050) = +46.2^\circ$, और
$2x - 1 = 46.2/66.0 = 0.700$: $x = 0.850$, \qty{85}{\percent}
$(+)$ \omterm{def:b1:stereochemistry-in-depth:enantiomers}{प्रतिबिंबरूपी} और \qty{15}{\percent} $(-)$।
\end{solution}

\begin{solution}{pb:b1:stereochemistry-in-depth:1}
\textbf{1.} साइक्लोहेक्सेन वलय: C1 पर OH, C2 पर आइसोप्रोपिल समूह, C5 पर
मेथिल; C1, C2 और C5 त्रिविमजनक हैं।
\textbf{2.} $2^3 = 8$, \omterm{def:b1:stereochemistry-in-depth:enantiomers}{प्रतिबिंबरूपियों} के चार युग्म।
\textbf{3.} किसी भी व्यवस्था में दर्पण-तल नहीं है: तीनों प्रतिस्थापी
सभी भिन्न हैं।
\textbf{4.} \ce{OH} $>$ C2 (C,C,H) $>$ C6 (C,H,H) $>$ H।
\textbf{5.} C1 (O,C,H) $>$ आइसोप्रोपिल (C,C,H) $>$ C3 (C,H,H) $>$ H।
\textbf{6.} C6 और C4 दोनों (C,H,H); आगे, C6, C1 (O,C,H) तक ले जाता है, C4, C3
(C,H,H) तक: C6 $>$ C4 $>$ \ce{CH3} $>$ H।
\textbf{7.} \cip{1S,2R,5S}।
\textbf{8.} \omterm{def:b1:stereochemistry-in-depth:conformations}{कुर्सी} में, सन्निकट कार्बनों (C1, C2) पर दो \omterm{def:b1:stereochemistry-in-depth:conformations}{विषुवतीय} आबंध
एक ऊपर और एक नीचे संकेत करते हैं: \textit{trans}; वलय के कार्बन 1 और 3 पर
(C1 और C5, C6 के माध्यम से) दोनों \omterm{def:b1:stereochemistry-in-depth:conformations}{विषुवतीय} आबंध एक ही ओर संकेत करते हैं:
\textit{cis}। यही मेंथॉल की व्यवस्था है।
\textbf{9.} तीनों \omterm{def:b1:stereochemistry-in-depth:conformations}{अक्षीय} हो जाते हैं।
\textbf{10.} पलटी हुई \omterm{def:b1:stereochemistry-in-depth:conformations}{कुर्सी} में अकेला मेथिल लगभग
\qty{5.7}{kJ/mol} की कीमत लेता और बड़ा आइसोप्रोपिल समूह अधिक; OH अपनी एक छोटी
कीमत जोड़ता है, और \omterm{def:b1:stereochemistry-in-depth:conformations}{अक्षीय} OH तथा मेथिल, कार्बन 1 और 3 पर और
एक ही फलक पर, एक-दूसरे पर भीड़ भी करते। यह \qty{12}{kJ/mol} से काफ़ी अधिक है,
अर्थात् $\exp(12000/2479) \approx 130$ से एक से अधिक का अनुपात: मेंथॉल
लगभग पूरी तरह पूर्णतः \omterm{def:b1:stereochemistry-in-depth:conformations}{विषुवतीय} \omterm{def:b1:stereochemistry-in-depth:conformations}{कुर्सी} में रहता है।
\textbf{11.} एक \omterm{def:b1:stereochemistry-in-depth:enantiomers}{अप्रतिबिंबरूपी}: तीन में से केवल एक केंद्र भिन्न है।
\textbf{12.} आइसोप्रोपिल और मेथिल \omterm{def:b1:stereochemistry-in-depth:conformations}{विषुवतीय} होने पर, OH \omterm{def:b1:stereochemistry-in-depth:conformations}{अक्षीय} है।
\textbf{13.} \omterm{def:b1:stereochemistry-in-depth:enantiomers}{अप्रतिबिंबरूपी} अपने सभी गुणों में भिन्न होते हैं (यहाँ
\omterm{def:b1:stereochemistry-in-depth:conformations}{अक्षीय} या \omterm{def:b1:stereochemistry-in-depth:conformations}{विषुवतीय} OH का हाइड्रोजन आबंधन); \omterm{def:b1:stereochemistry-in-depth:enantiomers}{प्रतिबिंबरूपी} केवल
\omterm{def:b1:stereochemistry-in-depth:stereocentre}{किरैल} साथियों के प्रति भिन्न होते हैं।
\textbf{14.} $[\alpha] = -2.46/(1.00 \times 0.0500) = -49.2^\circ$, संदर्भ-पुस्तिका के
परास के भीतर।
% ledger: alpha:l-menthol-range
\textbf{15.} $+49.2^\circ$; $0$।
\textbf{16.} $-1.97/0.0500 = -39.4^\circ$।
\textbf{17.} $\alpha = \ell c[\alpha]_{(-)}\bigl(x - (1 - x)\bigr) =
\ell c[\alpha]_{(-)}(2x - 1)$।
\textbf{18.} $x = \frac12(1 + \alpha/\alpha_{\text{संदर्भ}})$, समान $\ell$ और $c$ पर।
\textbf{19.} हाँ: कोई भी अन्य प्रकाशतः सक्रिय यौगिक अपना पद जोड़ता है
(\omterm{prop:b1:stereochemistry-in-depth:biot}{बायो का नियम})। घूर्णन को संघटन के रूप में पढ़ने से पहले किसी अन्य विधि
(वर्णलेखन) से रासायनिक शुद्धता की जाँच आवश्यक है।
\textbf{20.} $x = \frac12(1 + 1.97/2.46) = 0.900$।
\textbf{21.} $x = \frac12(1 + 2.40/2.46) = 0.988$।
\textbf{22.} दूसरा (\qty{98.8}{\percent})।
\textbf{23.} $-3.94^\circ$; तब वही पठन-त्रुटि कोण का छोटा
अंश होती है।
\textbf{24.} \textbf{\qty{90.0}{\percent} $(-)$-मेंथॉल} (और
\qty{10.0}{\percent} उसका \omterm{def:b1:stereochemistry-in-depth:enantiomers}{प्रतिबिंबरूपी})।
\end{solution}
